\documentclass[english, parskip=full]{scrreprt}
\usepackage[utf8]{inputenc}  % Kodierung der Datei
%\usepackage[ngerman]{babel}  % Sprache auf Deutsch 
\usepackage{color}
\usepackage{graphicx, gensymb}
\usepackage{amsfonts, cancel, amssymb}
\usepackage{amsmath, amsthm, array, esvect, esint}
\usepackage{booktabs, multirow}  % schönere Tabellen
\usepackage{caption}  % Anpassung der Captions
\usepackage{scrlayer-scrpage}    % Manipulation des Seitenstils
\pagestyle{scrheadings}  % Stil für die Seite setzen
\automark{section}  % Abschnittsnamen als Seitenbeschriftung 
\setheadsepline{.5pt}    % Kopzeile durch Linie abtrennen
\usepackage[hidelinks]{hyperref}  % Links und weitere PDF-Features
\usepackage{typearea, tikz, pgffor, verbatim}
\usetikzlibrary{calc, quotes}
\usepackage[cache=false, outputdir=build]{minted}

\definecolor{bg}{rgb}{0.9,0.9,0.9}
\newcommand\numberthis{\addtocounter{equation}{1}\tag{\theequation}}
\newcommand{\bra}{}
\newcommand{\kom}{}
\newcommand{\brak}{}
\newcommand{\braket}{}
\newcommand{\intx}{}
\newcommand{\intr}{}
\newcommand{\kla}{}
\newcommand{\klam}{}
\newcommand{\klammer}{}
\newcommand{\oper}{}
\newcommand{\tecxt}{}
\newcommand{\Bra}{}
\newcommand{\Ket}{}

\newcommand*{\titel}{Mathematical methods for Differential Equations}
\newcommand*{\autor}{Joachim Schwardt}

\hypersetup{pdfauthor={\autor}, pdftitle={\titel}}

\title{\titel}
\author{\autor}
\date{\today}

\begin{document}

\begin{titlepage}
\maketitle  % Titel setzen
\tableofcontents  % Inhaltsverzeichnis setzen
\end{titlepage}

\chapter{Introduction}
This book is a culmination of mathematical methods for solving various types of differential equations. Most examples are particularly important in physics, while some generalized ones are discussed purely for understanding their mathematical structure.
\chapter{Ordinary Differential Equations (ODE)}
General structure...
\section{Linear homogeneous ODE}
Subclass definition...
\subsection{Constant coefficients}
Definition of constant coefficients...
\subsubsection{First order}\label{sss:ODE linhom cc first order}
our first ode...
\subsubsection{Second order}
harmonic oscillator...
\subsubsection{Higher orders}
general solution...
\section{Linear inhomogeneous ODE}
much more stuff...
\chapter{Partial Differential Equations (PDE)}
\section{Linear homogeneous PDE}
Let $f:\mathbb{R}^N\rightarrow\mathbb{R}$
\subsection{Constant coefficients}
Linear PDEs with constant coefficients can be solved by Fourier transformations:
\begin{align*}
\mathcal{F}_N[f](k)&:=(2\pi)^{-N/2}\intr\int_{\mathbb{R}^{N}}{f(x)e^{-ik\cdot x}}\,\mathrm{d}^{N}x \\
\mathcal{F}_N^{-1}[f](x)&:=(2\pi)^{-N/2}\intr\int_{\mathbb{R}^{N}}{f(k)e^{ik\cdot x}}\,\mathrm{d}^{N}k
\end{align*}
Note that we will use the notation $f(k):=\mathcal{F}[f](k)$ for any function $f(x)$. This is ambiguous, since $f(k)$ and $f(x)$ do \textsl{not} refer to the same function. We only differentiate between the function $f(x)$ in direct space and $f(k)$ in reciprocal space by their arguments.
\subsubsection{First order}
Let $x^\mu,a^\mu\in\mathbb{R}^{N+1}$:
\begin{align*}
0&=f(x^\mu) + a^\mu\partial_\mu f(x^\mu)
\end{align*}
Consider the $N$-dimensional Fourier transformation, leaving us with the arguments $t,k$
\begin{align*}
0&=\mathcal{F}_{N}[f+a\cdot\nabla f](t,k)=f(t,k)+a_0\partial_{t}f(t,k) + ia\cdot k  f(t,k)
\end{align*}
Now this is just an ODE! Consider the $k$ as parameters and apply our knowledge from (\ref{sss:ODE linhom cc first order}). Note that the arbitrary constant is now a function of $k$ and define the ``frequency'' $\Omega:=\frac{1+ia\cdot k}{a_0}$:
\begin{align*}
f(t,k)&=A(k)e^{\Omega t}
\end{align*}
Our initial condition is given by $\phi(t=0,x)=:\phi(x)$. Using this, we can derive an expression for $A(k)$:
\begin{align*}
A(k)&=f(0,k)=\mathcal{F}_{N}[\phi](k)=(2\pi)^{-N/2}\intr\int_{\mathbb{R}^{N}}{\phi(x)e^{-ik\cdot x}}\,\mathrm{d}^{N}x
\end{align*}
With this, we can solve the PDE:
\begin{align*}
f(x^\mu)&=\mathcal{F}_{N}^{-1}[f(t,k)](x)=\mathcal{F}_N^{-1}\klam\left[ {e^{\Omega t}\mathcal{F}_N[\phi]} \right](t,x)=\\
&=(2\pi)^{-N}\intr\int_{\mathbb{R}^{N}}{e^{\frac{1+ia\cdot k}{a_0}t}\intr\int_{\mathbb{R}^{N}}{\phi(x)e^{-ik\cdot x}}\,\mathrm{d}^{N}x}\,e^{ik\cdot x}\,\mathrm{d}^{N}k
\end{align*}
We want to discuss this solution in a bit more detail by considering a few examples. Let $\phi\equiv\delta$:
\begin{align*}
f(x^\mu)&=(2\pi)^{-N}\intr\int_{\mathbb{R}^{N}}{e^{\frac{1+ia\cdot k}{a_0}t+ik\cdot x}}\,\mathrm{d}^{N}k=e^{t/a_0}\prod_{j=1}^N\frac{1}{2\pi}\intr\int_{\mathbb{R}^{ }}e^{ik_j(a_jt/a_0 + x_j)}\,\mathrm{d}^{}k_j=\\
&=e^{t/a_0}\prod_{j=1}^N\delta\kla\left( {x_j+\frac{a_jt}{a_0}} \right)=\delta\kla\left( {x+\frac{at}{a_0}} \right)e^{t/a_0}
\end{align*}
\subsubsection{Second order}
While technically of second order, the heat equation can be solved by directly applying the method from above, because it is still only of first order in one variable:
\begin{align*}
0&=\partial_tf(t,x)-D\nabla^2f(t,x)
\end{align*}
Apply the $N$-dimensional Fourier transformation:
\begin{align*}
0&=\mathcal{F}[\partial_tf-D\nabla^2f](t,k)=\partial_tf(t,k)+Dk^2f(t,k)
\end{align*}
Once again solve the resulting ODE using (\ref{sss:ODE linhom cc first order}):
\begin{align*}
f(t,k)&=A(k)e^{-Dk^2t}
\end{align*}
Find $A(k)$ from the initial condition $f(t=0,x)=:\phi(x)$:
\begin{align*}
A(k)&=f(0,k)=\mathcal{F}_N[\phi](k)=(2\pi)^{-N/2}\intr\int_{\mathbb{R}^{N}}{\phi(x)e^{-ik\cdot x}}\,\mathrm{d}^{N}x
\end{align*}
And solve for $f$ by inverse Fourier transformation:
\begin{align*}
f(t,x)&=\mathcal{F}_N^{-1}\klam\left[ {e^{-Dk^2t}\mathcal{F}_N[\phi]} \right]=\\
&=(2\pi)^{-N}\intr\int_{\mathbb{R}^{N}}{e^{-Dk^2t}\intr\int_{\mathbb{R}^{N}}{\phi(x)e^{-ik\cdot x}}\,\mathrm{d}^{N}x}\,e^{ik\cdot x}\,\mathrm{d}^{N}k
\end{align*}
Now, just as above, we can not simplify this further without additional knowledge about $\phi$. However, for the special case of $\phi\equiv\delta$, we arrive at the characteristic solution of a bell curve with increasing width:
\begin{align*}
f(t,x)&=(2\pi)^{-N}\intr\int_{\mathbb{R}^{N}}{e^{-Dk^2t+ik\cdot x}}\,\mathrm{d}^{N}k=(4\pi Dt)^{-N/2}e^{-\frac{x^2}{4Dt}}
\end{align*}
Next, we will approach ``true'' second order PDE, namely the Laplace equation:
\begin{align*}
0&=\nabla^2f(x^\mu)
\end{align*}
Fourier transformation and solution of the resulting ODE yield
\begin{align*}
0&=\mathcal{F}_N[\nabla^2f](t,k)=\partial_t^2f(t,k)-k^2f(t,k)\\
f(t,k)&=A(k)e^{kt}+B(k)e^{-kt}\\
\partial_tf(t,k)&=kA(k)e^{kt}-kB(k)e^{-kt}
\end{align*}
Given the initial conditions $f(t=0,x)=:\phi(x)$ and $\partial_tf(t=0,x)=:\psi(x)$, we can determine $A(k)$ and $B(k)$:
\begin{align*}
2A(k)&=f(0,k)+\frac{1}{k}\partial_tf(0,k)=\mathcal{F}_N[\phi](k)+\frac{1}{k}\mathcal{F}_N[\psi](k)\\
2B(k)&=f(0,k)-\frac{1}{k}\partial_tf(0,k)=\mathcal{F}_N[\phi](k)-\frac{1}{k}\mathcal{F}_N[\psi](k)
\end{align*}
From this, we can extract our solution:
\begin{align*}
f(t,x)&=\mathcal{F}_N^{-1}\klam\left[ {A(k)e^{kt}+B(k)e^{-kt}} \right]
\end{align*}
\subsubsection{Higher orders}
We can generalize this approach for an arbitrary linear homogeneous PDE with constant coefficients. For this, define a general differential operator accordingly:
\begin{align*}
L&:=
\end{align*}
%\begin{minted}[frame=lines, bgcolor=bg, fontsize=\footnotesize, linenos]{python}

\begin{thebibliography}{9}
\bibitem{example} description, \url{https://www.exmapleurl}, day.\,month~2021.
\end{thebibliography}

\appendix
\chapter{Some more details}

\end{document}